% !TEX root = ../Main.tex

% === TexCount-Definitions ===
% (Ignore Headings)
%TC:macro \chapter [ignore]
%TC:macro \section [ignore]
%TC:macro \subsection [ignore]
%TC:macro \subsubsection [ignore]
%TC:macro \includesvg [ignore]
% (Ignore Figures and Tables)
%TC:envir figure [ignore] ignore
%TC:envir table [ignore] ignore

% === Content ===
\newpage
% !TEX root = ../../Main.tex
% ==============================================================================
% Kapitel 1: Einleitung
% ==============================================================================
% Description: Einleitung zur Formatierung der Seminararbeit
% ==============================================================================

\chapter{Einleitung}
\label{chap:einleitung}

Das hier vorliegende Dokument soll als Vorlage für die Gestaltung der im Seminarkurs \emph{Ausgewählte Themen der Medizinisch-Technischen Informatik} erstellten schriftlichen Ausarbeitungen dienen.

Es ist zu beachten, dass das Formatieren des Textes als einer der letzten Schritte der Ausarbeitung durchgeführt wird, da dieser Schritt erfahrungsgemäß viel Zeit in Anspruch nimmt und daher nur einmalig ausgeführt werden sollte.

Der vorliegende Entwurf wurde mit \LaTeX{} unter Verwendung der HSRTReport-Dokumentklasse erstellt, welche auf der KOMA-Script-Klasse \texttt{scrreprt} basiert. Diese Vorlage bietet eine professionelle und konsistente Formatierung für wissenschaftliche Arbeiten.

Alle in diesem Dokument enthaltenen Hinweise zur Gestaltung des Dokumentes dienen als Referenz für die korrekte Verwendung des Templates. Die Vorlage selbst übernimmt bereits die meisten Formatierungsanforderungen automatisch.

% ==============================================================================
% End of Kapitel 1
% ==============================================================================

%!TEX root = ../../main.tex
% ==============================================================================
% CHAPTER 2: WISSENSCHAFTLICHE EVIDENZ
% ==============================================================================
\chapter{Wissenschaftliche Evidenz für Enhancement durch Hirnstimulation}
\label{chap:evidenz}

% LITERATURVORSCHLÄGE FÜR KAPITEL:
% - Simonsmeier et al. (2018): "Electrical brain stimulation improves learning" [Meta-Analyse]
% - Reis et al. (2009): "Noninvasive cortical stimulation enhances motor skill acquisition" [motorisch, 1700+ Zit.]
% - Akkad et al. (2021): "Increasing human motor skill acquisition by driving theta-gamma" [innovative Studie]
% - Cavaleiro et al. (2020): "Memory and Cognition-Related Neuroplasticity" [Mechanismen]
% - Frontiers meta-analysis (2024): TMS effects on cognition
% - Vergallito et al. (2022): "Inter-Individual Variability in tDCS Effects" [Variabilität, 154 Zit.]

\section{Neurobiologische Wirkmechanismen}
\label{sec:wirkmechanismen}

% LITERATUR: Cavaleiro 2020, Woods 2016 technical guide, Niksche & Paulus reviews

% \subsection{Grundprinzipien der transkraniellen Hirnstimulation}
% \label{subsec:grundprinzipien}
% - Elektromagnetische Induktion
% - Modulation kortikaler Erregbarkeit
% - Unterschiede tDCS (konstant), tACS (oszillatorisch), TMS (magnetisch)
% - Relevante Hirnregionen (DLPFC, M1, parietal cortex)

% Subsubsection 2.1.1.1: tDCS Mechanismus
% \subsubsection{Transkranielle Gleichstromstimulation (tDCS)}
% \label{subsubsec:tdcs_mech}
% - Anodale vs. kathodale Stimulation
% - Membranpotenzial-Verschiebung
% - Stromfluss und elektrische Feldverteilung
% - Parameter: Intensität (1-2 mA), Dauer (20-30 min), Elektroden-Montage

% Subsubsection 2.1.1.2: tACS Mechanismus
% \subsubsection{Transkranielle Wechselstromstimulation (tACS)}
% \label{subsubsec:tacs_mech}
% - Frequenzabhängige Modulation
% - Entrainment neuronaler Oszillationen
% - Theta (4-8 Hz), Gamma (40 Hz), andere Bänder
% - Theta-Gamma-Kopplung für Gedächtnis/Lernen

% Subsubsection 2.1.1.3: TMS/rTMS Mechanismus
% \subsubsection{Transkranielle Magnetstimulation (TMS/rTMS)}
% \label{subsubsec:tms_mech}
% - Faradaysche Induktion
% - Aktionspotenziale in kortikalen Neuronen
% - High-frequency (>5 Hz) exzitatorisch, Low-frequency (<1 Hz) inhibitorisch
% - iTBS vs. cTBS Protokolle

% \subsection{Zelluläre und molekulare Plastizitätsmechanismen}
% \label{subsec:plastizität}

% LITERATUR: Cavaleiro 2020, Frontiers molecular neuroscience 2020, Esser 2006

% Subsubsection 2.1.2.1: Long-Term Potentiation (LTP)
% \subsubsection{Langzeitpotenzierung und -depression}
% \label{subsubsec:ltp_ltd}
% - LTP als synaptische Verstärkung
% - LTD als synaptische Abschwächung
% - iTBS induziert LTP-ähnliche Effekte
% - cTBS induziert LTD-ähnliche Effekte
% - Messung: MEP (Motor Evoked Potentials)

% Subsubsection 2.1.2.2: Neurotrophische Faktoren
% \subsubsection{Neurotrophische Faktoren und Genexpression}
% \label{subsubsec:neurotrophisch}
% - BDNF (Brain-Derived Neurotrophic Factor)
% - c-fos, zif268 (immediate early genes)
% - NMDA-Rezeptoren
% - Hochregulation durch Stimulation (besonders iTBS)

% Subsubsection 2.1.2.3: Transmitter-Modulation
% \subsubsection{Neurotransmitter-Modulation}
% \label{subsubsec:transmitter}
% - GABAerge und glutamaterge Systeme
% - Dopamin und Lernverhalten (Reward-based Learning)
% - Monoamine (Serotonin, Noradrenalin)
% - Inter-individuelle Variabilität durch Transmitter-Polymorphismen

% \subsection{Zeitliche Aspekte: Phasen des Lernens}
% \label{subsec:zeitliche_aspekte}

% LITERATUR: Reis 2009, Debarnot 2019, Zimerman 2013

% Subsubsection 2.1.3.1: Enkodierung und Online-Lernen
% \subsubsection{Enkodierung und Online-Lernen}
% \label{subsubsec:enkodierung}
% - Effekte während aktiver Aufgabe
% - Stimulation während Trainingsphase
% - Hemmt oder fördert Online-Lernen je nach Timing/Polarität
% - Studien: [Refs: Reis 2009 - motorisch, Stagg 2011]

% Subsubsection 2.1.3.2: Konsolidierung und Offline-Effekte
% \subsubsection{Konsolidierung und Offline-Effekte}
\label{subsubsec:konsolidierung}
% - Verstärkung zwischen Trainings-Sessions
% - Sleep-dependent consolidation
% - tDCS während Tiefschlaf: Slow Oscillations induzieren
% - Externe Effekte: 1h, 24h, 1 Woche post-stimulation
% - Studien: [Reis 2009, Marshall & Born 2007]

\section{Empirische Befunde nach kognitiver Domäne}
\label{sec:befunde_domänen}

% LITERATURVORSCHLÄGE: Simonsmeier 2018, Reis 2009, Akkad 2021, Lopez-Alonso 2025

% \subsection{Motorisches Lernen und Fertigkeitserwerb}
\label{subsec:motorisch}

% LITERATUR: Reis 2009 (1700 Zit.), Akkad 2021, Maceira-Elvira 2024, Lopez-Alonso 2025

% Subsubsection 2.2.1.1: Anodales tDCS über M1
% \subsubsection{Anodale tDCS über motorischem Kortex (M1)}
\label{subsubsec:motor_tdcs}
% - Effekte auf Konsolidierung (nicht Online-Lernen)
% - Protokoll: 2 mA, 20 min, posttraining
% - Effektgrößen: Cohen's d = 0.5-1.0
% - Persistenz: bis zu 24h
% - Alter-abhängig: auch ältere Probanden profitieren

% Subsubsection 2.2.1.2: Theta-Gamma-tACS über M1/sensomotorischer Kortex
% \subsubsection{Theta-Gamma-tACS für motorisches Lernen}
\label{subsubsec:motor_tacs}
% - Innovative Frequenzkopplung
% - Phase-Amplitude-Coupling (4-8 Hz Theta moduliert 40 Hz Gamma)
% - Replikationsstudie [Akkad 2021]
% - Effekte: Beschleunigung, Persistenz über 1h
% - Mechanismus: Theta-Gamma relevant für hippokampales Lernen

% Subsubsection 2.2.1.3: Komplexe motorische Skills
% \subsubsection{Komplexe motorische Skills (negative Befunde)}
\label{subsubsec:komplexe_motor}
% - Golf-Putting Studie [Lopez-Alonso 2025]
% - tDCS zeigte KEINE zusätzliche Verbesserung über Training
% - Implikation: Einfache vs. komplexe Skills unterschiedlich
% - Limitationen: Realwelt-Komplexität, Multi-Link Bewegungen

% \subsection{Akademische Fähigkeiten (Mathematik, Sprache)}
\label{subsec:akademisch}

% LITERATUR: Simonsmeier et al. 2018 (133 Zit., Meta-Analyse), Cohen Kadosh Arbeiten, Sarkar 2014

% Subsubsection 2.2.2.1: Timing der Stimulation
% \subsubsection{Timing: Stimulation während Lernphase vs. Test}
\label{subsubsec:timing_learning}
% - Simonsmeier Meta-Analyse [2018]: Stimulation während Lernen > während Test
% - Mechanismus: LTP-Induktion während aktiver Verarbeitung
% - Effektgrößen: SMD = 0.3-0.7 für akademische Fähigkeiten

% Subsubsection 2.2.2.2: Mathematische Kompetenzen
% \subsubsection{Mathematische Fähigkeiten}
\label{subsubsec:mathe}
% - tRNS (Random Noise) über "Mathezentrum" (DLPFC/IPS)
% - Schmerzfreie Alternative zu tDCS
% - Verbesserung Reaktionszeit + Genauigkeit
% - [Cohen Kadosh Arbeiten, erwähnt in Spektrum.de]

% Subsubsection 2.2.2.3: Sprachfähigkeiten
% \subsubsection{Sprachfähigkeiten}
\label{subsubsec:sprache}
% - Weniger Forschung als Mathematik
% - tDCS über Sprach-Arealen (Broca, Wernicke)
% - Stimulation während Vokabel-Lernen
% - Effektgrößen: moderate bis klein

% \subsection{Arbeitsgedächtnis und exekutive Funktionen}
\label{subsec:workingmem}

% LITERATUR: Hoy et al. 2016, Senkowski & Sobirey 2022, Imburgio & Orr 2018

% Subsubsection 2.2.3.1: tDCS vs. tACS Vergleich
% \subsubsection{Vergleich tDCS vs. tACS auf Arbeitsgedächtnis}
\label{subsubsec:tdcs_vs_tacs}
% - Systematische Review [Senkowski 2022]: 43 Studien, 1826 Teilnehmer
% - Single-Session tDCS: kaum signifikant
% - Multi-Session tDCS: moderate Effekte (SMD ~0.3)
% - tACS (v.a. Gamma): robustere Effekte (SMD ~0.4-0.6)
% - Frequenz-Abhängigkeit: Gamma (40 Hz) > andere

% Subsubsection 2.2.3.2: iTBS auf DLPFC
% \subsubsection{Intermittierende Theta-Burst-Stimulation (iTBS) auf DLPFC}
\label{subsubsec:itbs_dlpfc}
% - Hoy et al. (2016): iTBS > tDCS auf 3-back Working Memory Task
% - Theta & Gamma Synchronisation erhöht
% - Effektgrößen größer als tDCS
% - Schnellere Behandlung (10min statt 20min)

% Subsubsection 2.2.3.3: Exekutive Funktionen
% \subsubsection{Exekutive Funktionen (Inhibition, Flexibilität)}
\label{subsubsec:exek_funkt}
% - tDCS Effekte: limitiert auf Refresh (WM)
% - Inhibition: kleine bis keine Effekte [Imburgio & Orr 2018]
% - Kognitive Flexibilität: wenig untersucht
% - Kontext-abhängigkeit der Effekte

% \subsection{Praktische Anwendungen in der Realwelt}
\label{subsec:realwelt_apps}

% LITERATUR: Clark et al. 2012, Lopez-Alonso 2025, Violante 2024

% Subsubsection 2.2.4.1: Threat-Detection-Training (Militär)
% \subsubsection{Threat-Detection-Training}
\label{subsubsec:threat_detection}
% - Clark et al. (2012): Virtual Reality + tDCS
% - 96 Probanden, Erkennung versteckter Objekte
% - Signifikante Verbesserung + Trainingstransfer
% - DARPA finanzierte Forschung
% - Praktische Implikation: funktioniert in komplexer Umgebung

% Subsubsection 2.2.4.2: Negative Befunde und Realwelt-Limitationen
% \subsubsection{Negative Befunde und Realwelt-Limitationen}
\label{subsubsec:negative_realwelt}
% - Lopez-Alonso (2025): Golf-Putting keine zusätzliche Verbesserung
% - Violante et al. (2024): "modest evidence" für Neuro-Enhancement
% - Komplexe Skills brauchen Integration verschiedener Systeme
% - Transfer-Problem: Lab ≠ Realwelt

\section{Kritische Limitationen der Evidenz}
\label{sec:limitationen}

% LITERATURVORSCHLÄGE: Vergallito 2022, Chew 2015, Maceira-Elvira 2024, Violante 2024

% \subsection{Inter-individuelle Variabilität in Stimulationsresponse}
\label{subsec:variabilität}

% LITERATUR: Vergallito et al. 2022 (154 Zit.), Chew et al. 2015 (317 Zit.), Filmer et al. 2014

% Subsubsection 2.3.1.1: Anatomische Unterschiede
% \subsubsection{Anatomische Unterschiede}
\label{subsubsec:anat_diff}
% - Kortikale Gyrierung (Gyri/Sulci Patterns)
% - Kortexdicke in spezifischen Regionen
% - Faserverbindungen zwischen Regionen
% - Vorhersage: Dicke in mittlerem Stirnfurchen prädiziert tDCS-Response [Filmer]
% - Modellierung: FEM (Finite-Elemente) kann Feldverteilung individualisieren

% Subsubsection 2.3.1.2: Genetische und biochemische Faktoren
% \subsubsection{Genetische und biochemische Faktoren}
\label{subsubsec:genetik}
% - BDNF val66met Polymorphismus
% - Transmitter-Polymorphismen (COMT, DAT)
% - Baseline-GABA/Glutamat Ratio
% - Limitierte Forschung, aber vielversprechend für Personalisierung

% Subsubsection 2.3.1.3: Funktionelle Baseline-Unterschiede
% \subsubsection{Funktionelle Baseline-Unterschiede}
\label{subsubsec:baseline_func}
% - Resting-State Konnektivität
% - Theta-Power im Frontallappen (prädiziert rTMS-Response bei Depression)
% - Baseline-Leistung auf kognitiver Aufgabe

% \subsection{Baseline-Leistung und Lerner-Typ}
\label{subsec:baseline_leistung}

% LITERATUR: Maceira-Elvira 2024, Sarkar 2014

% Subsubsection 2.3.2.1: Wer profitiert am meisten?
% \subsubsection{Baseline-abhängige Responsiveness}
\label{subsubsec:baseline_response}
% - Maceira-Elvira (2024): "Native learning ability, not age, determines response"
% - Personen mit niedrigen Ausgangsfähigkeiten profitieren MEHR
% - Hochleister: KEINE zusätzliche Verbesserung durch tDCS
% - Implikation: Enhancement-Effekt abhängig vom Starting Point

% Subsubsection 2.3.2.2: Math-Anxiety Beispiel
% \subsubsection{Kontext-abhängige Effekte: Math-Anxiety}
\label{subsubsec:math_anxiety}
% - Sarkar et al. (2014): tDCS hilft bei Mathe-Angst
%   - Verbesserung Reaktionszeit + reduziertes Cortisol
% - Aber: Bei Nicht-Angst-Personen verschlechtert sich Leistung!
% - Implikation: Effekte nicht universal, stark kontext-abhängig

% \subsection{Nur 39-45\% sind Responder}
\label{subsec:responder_rate}

% LITERATUR: Vergallito 2022, Chew 2015, Filmer 2014

% Subsubsection 2.3.3.1: Empirische Evidenz
% \subsubsection{Responder vs. Non-Responder Phänomenologie}
\label{subsubsec:responder_phenom}
% - Vergallito (2022): Nur 39-45% zeigen erwartete Effekte
% - Cluster-Analysen zeigen bimodale Verteilungen
% - Chew (2015): Intra-individual Variability ICC = -0.50 (negligible!)
% - Problem: Nicht zuverlässig vorhersagbar mit Current Protocols

% Subsubsection 2.3.3.2: Implikationen für praktische Anwendung
% \subsubsection{Implikationen für praktische Anwendung}
\label{subsubsec:responder_impl}
% - Screening notwendig zur Identification von Respondern
% - Adaptives Protokoll-Adjustments
% - Oder: Personalisierung VOR Behandlung (individuelles FEM-Modell, fMRI)

% \subsection{Transferdefizite und Aufgabenspezifität}
\label{subsec:transfer}

% LITERATUR: Violante 2024, Lopez-Alonso 2025

% Subsubsection 2.3.4.1: Begrenzte Generalisierbarkeit
% \subsubsection{Begrenzte Generalisierbarkeit der Effekte}
\label{subsubsec:generalisierbar}
% - Effekte meist aufgabenspezifisch
% - Transfer zu anderen Aufgaben limitiert
% - Transfer zu anderen Domänen: unklar
% - Implikation: Nicht "allgemeine Intelligenz-Verbesserung"

% Subsubsection 2.3.4.2: Realwelt-Transfer
% \subsubsection{Transfer von Lab zu Realwelt}
\label{subsubsec:lab_realwelt}
% - Threat-Detection erfolgreich [Clark 2012]
% - Golf-Putting KEIN Transfer [Lopez-Alonso 2025]
% - Komplexität des Realwelt-Skills kritisch
% - Context-Dependence der Stimulation-Effekte

% \subsection{Fehlende Langzeitstudien und Persistenz}
\label{subsec:langzeit}

% LITERATUR: Brühl 2019, Woodham 2024 (nur 10 Wochen), Violante 2024

% Subsubsection 2.3.5.1: Kurzfristige vs. Langfristige Effekte
% \subsubsection{Persistenz von Effekten}
\label{subsubsec:persistenz}
% - Kurzfristig: Gut dokumentiert (bis 24h post-stimulation)
% - Mittelfristig (1-4 Wochen): Noch nicht vollständig untersucht
% - Langfristig (Monate/Jahre): Kaum Daten
% - Frage: Was passiert nach Beendigung der Stimulation?

% Subsubsection 2.3.5.2: Lifetime Safety bei Gesunden
% \subsubsection{Lifetime Safety bei wiederholter Anwendung bei Gesunden}
\label{subsubsec:lifetime_safety}
% - Kurzfristige Nebenwirkungen: gut dokumentiert, mild
% - Langfristige Effekte unbekannt [Brühl 2019]
% - Besondere Bedenken: Hirnentwicklung (<25 Jahre)
% - Frage: Könnte wiederholte Stimulation Neuroplastizität negativ beeinflussen?
% - Forschungsbedarf: Prospektive Langzeitstudien

%!TEX root = ../../main.tex
% ==============================================================================
% CHAPTER 3: PRAKTIKABILITÄT
% ==============================================================================
\chapter{Praktikabilität in Bildungs- und Trainingskontexten}
\label{chap:praktikabilität}

% LITERATURVORSCHLÄGE FÜR KAPITEL:
% - Cappon et al. (2023): "Home-based tES training program" [Machbarkeit, 10 Zit.]
% - Woodham et al. (2024): "Home-based tDCS in MDD" [Nature Medicine, 62 Zit., praktisches Modell]
% - Woods et al. (2015/2016): "Technical guide to tDCS" [Standardwerk, 1700+ Zit.]
% - Thair et al. (2017): "Transcranial Direct Current Stimulation (tDCS)" [Praktische Details]
% - Caulfield et al. (2022): "Optimized electrode positions" [Technische Optimierung]

\section{Technische Aspekte und Durchführbarkeit}
\label{sec:tech_durchführbarkeit}

% LITERATUR: Woods 2015/2016, Thair 2017, Caulfield 2022

% \subsection{Kosteneffizienz und ökonomische Zugänglichkeit}
\label{subsec:kosten}

% LITERATUR: neurocare academy, Woodham 2024, Cappon 2023

% Subsubsection 3.1.1.1: Gerätekosten
% \subsubsection{Gerätekosten und Verfügbarkeit}
\label{subsubsec:gerätekosten}
% - tDCS Consumer-Geräte: $200-500
% - rTMS klinische Geräte: $10,000-50,000
% - tACS/tRNS: mittleres Preissegment
% - Vergleich:
%   - Private Nachhilfe/Einzelunterricht: €30-60/h
%   - Elite-Trainingsprogramme: €1000+/Monat
%   - tDCS: ~€5-10 pro Sitzung (Amortisation)

% Subsubsection 3.1.1.2: Kostenvergleich zu Alternativen
% \subsubsection{Kostenvergleich zu alternativen Interventionen}
\label{subsubsec:kostenvergleich}
% - Vs. Einzelunterricht: tDCS kostengünstiger bei Skalierung
% - Vs. Medikamente (Modafinil, etc.): vergleichbar
% - Vs. kognitive Trainings-Apps: tDCS höher initial, aber kombinierbar

% \subsection{Heimbasierte und dezentrale Anwendung}
\label{subsec:heimbasiert}

% LITERATUR: Cappon 2023 (10 Zit., Pilotprogramm), Woodham 2024 (Nature Med, 62 Zit., 174 Probanden)

% Subsubsection 3.1.2.1: Erfolgreiche Pilotprojekte
% \subsubsection{Erfolgreiche Pilotprojekte mit Heimanwendung}
\label{subsubsec:pilotprojekte}
% - Cappon et al. (2023): Home-based tES Trainings-Programm
%   - Laien-Administration nach Trainig
%   - Adhärenz: 98-100%
%   - Sicherheit demonstriert
%   - Downloadbare Materialien verfügbar
% - Woodham et al. (2024, Nature Medicine): Home-based tDCS bei MDD
%   - 174 Teilnehmer, remote supervision
%   - 10 Wochen Protokoll
%   - Hohe Akzeptanz, gute Effektivität

% Subsubsection 3.1.2.2: Remote Supervision und Monitoring
% \subsubsection{Remote Supervision und Monitoring}
\label{subsubsec:remote_supervision}
% - Möglichkeit mit modernen Technologien
% - Telemedizin-Modelle existieren
% - Dokumentation via Apps/Cloud
% - Sicherheit: Contact person im Notfall nötig

% Subsubsection 3.1.2.3: Skalierbarkeit für breite Populationen
% \subsubsection{Skalierbarkeit für breite Populationen}
\label{subsubsec:skalierbarkeit}
% - Implication: Theoretisch skalierbar
% - Unterschied:
%   - Élite-Training (kleine Gruppen, hohe Personalkosten): möglich
%   - Mainstream-Bildung (Millionen Schüler): logistisch schwierig
%   - Hybrid-Modell: vielversprechend (Schulungen + Home-Use)

% \subsection{Sicherheit und Nebenwirkungsprofil}
\label{subsec:sicherheit}

% LITERATUR: Woods 2016 (technischer Standard), Wassermann-Kriterien, McCambridge et al. Sicherheit-Reviews

% Subsubsection 3.1.3.1: Häufige milde Nebenwirkungen
% \subsubsection{Häufige und milde Nebenwirkungen}
\label{subsubsec:mild_nebenwirkungen}
% - Kopfschmerzen: 10-20% (mit rezeptfreien Schmerzmitteln behandelbar)
% - Kopfhaut-Irritation/Brennen: lokal, vorübergehend
% - Muskelzuckungen (bei TMS): während Stimulation, gewöhnung nach wenigen Sessions
% - Müdigkeit/Benommenheit: selten, vorübergehend (<1h)
% - Übelkeit: sehr selten

% Subsubsection 3.1.3.2: Schwerwiegende Komplikationen
% \subsubsection{Schwerwiegende Komplikationen}
\label{subsubsec:severe_complications}
% - Epileptische Anfälle: extrem selten (<0.1% bei compliance mit Wassermann-Kriterien)
% - Keine bekannten Fälle permanenter Hirnschädigung
% - Vergleich zu Medikamenten: günstiger Effekt-Nebenwirkung-Profile
% - Kontraindikationen gut etabliert

% Subsubsection 3.1.3.3: Vergleich zu pharmakologischen Enhancern
% \subsubsection{Vergleich zu pharmakologischen Enhancern}
\label{subsubsec:pharm_vs_stim}
% - tDCS: keine systemischen Effekte (lokal wirksam)
% - Modafinil/Amphetamine: systemische Nebeneffekte (kardiovaskular, GI, etc.)
% - Koffein: Entzugssymptome, Toleranz
% - Fazit: tDCS sicherere Alternative aus Nebenwirkungs-Perspektive

\section{Praktische Implementierungshürden}
\label{sec:hürden}

% LITERATURVORSCHLÄGE: Vergallito 2022, Woods 2016, Meinzer 2024, Cappon 2023

% \subsection{Personalisierung und Vorhersage von Responsiveness}
\label{subsec:personalisierung}

% LITERATUR: Vergallito 2022, Meinzer 2024 (neuroimaging context), computational modeling papers

% Subsubsection 3.2.1.1: Das Problem: 60\% Non-Responder
% \subsubsection{Das Problem der 60\% Non-Responder}
\label{subsubsec:nonresponder_problem}
% - Vergallito (2022): Nur 39-45% zeigen erwartete Effekte
% - Current Standardprotokolle: "One-Size-Fits-All" funktioniert nicht
% - Screening nötig ODER Individualisierung VOR Behandlung

% Subsubsection 3.2.1.2: Anatomische Personalisierung
% \subsubsection{Anatomische Personalisierung}
\label{subsubsec:anat_personalisierung}
% - MRT-basierte Modelle: Individuelle Elektrodenplatzierung
% - Finite-Elemente-Modelle (FEM): Stromfluss berechnen
% - Toolboxen: SimNIBS, COMETS (MATLAB)
% - Kosten/Nutzen: Trade-off zwischen Präzision und praktische Durchführbarkeit

% Subsubsection 3.2.1.3: Funktionelle und genetische Personalisierung
% \subsubsection{Funktionelle und genetische Biomarker}
\label{subsubsec:biomarker_personal}
% - Theta-Power im Frontallappen (prädiziert rTMS-Response)
% - BDNF-Polymorphismen, COMT, DAT
% - EEG-basierte adaptive Protokolle (Closed-Loop)
% - Forschungsbedarf: Validierung dieser Biomarker

% \subsection{Standardisierung und Schulung von Anwendern}
\label{subsec:standardisierung}

% LITERATUR: Cappon 2023 (Trainings-Programm), neurocare academy, IFCN Guidelines

% Subsubsection 3.2.2.1: Notwendigkeit von Standardisierten Protokollen
% \subsubsection{Standardisierte Protokolle und Guidelines}
\label{subsubsec:standardisierte_protokolle}
% - Unterschiede aktuelle Praxis: große Variabilität
% - IFCN Guidelines (International Federation of Clinical Neurophysiology)
% - Wassermann-Kriterien für Sicherheit
% - Notwendig: Einigung auf "Best Practice" Protokolle

% Subsubsection 3.2.2.2: Trainingsmaterialien für nicht-medizinisches Personal
% \subsubsection{Trainingsmaterialien für Laien}
\label{subsubsec:training_laien}
% - Cappon (2023) hat Programm entwickelt: Video, Quizze, Checklisten
% - Erfolgreiche Durchführbarkeit demonstriert
% - Unterschied: Laien können Geräte anlegen, aber müssen medizinisch überwacht werden

% Subsubsection 3.2.2.3: Regulatorische Fragen zur Anwender-Qualifikation
% \subsubsection{Regulatorische Fragen}
\label{subsubsec:regulatory_qual}
% - Wer darf tDCS/TMS anwenden?
% - Ärzte, Therapeuten, Trainer, Laien?
% - Versicherungsdeckung, Haftung
% - Zertifizierungspfade nötig

% \subsection{Realistische Effektgrößen und Erwartungsmanagement}
\label{subsec:effektgrößen}

% LITERATUR: Reis 2009, Akkad 2021, Simonsmeier 2018, Violante 2024

% Subsubsection 3.2.3.1: Evidenzbasierte Effektgrößen
% \subsubsection{Empirische Effektgrößen aus Studien}
\label{subsubsec:empirisch_effekt}
% - Motorisches Lernen: Beschleunigung ~10-30% [Reis 2009, Akkad 2021]
% - Akademische Fähigkeiten: SMD = 0.3-0.7 [Simonsmeier 2018]
% - Arbeitsgedächtnis (Multi-Session): SMD ~0.3-0.5
% - NICHT: Keine "Wunder-Effekte"

% Subsubsection 3.2.3.2: Vergleich zu anderen Lerntechniken
% \subsubsection{Benchmarking gegen etablierte Methoden}
\label{subsubsec:benchmarking}
% - Spaced Repetition: SMD ~0.5-1.0
% - Elaboration: SMD ~0.7
% - tDCS: SMD ~0.3-0.7 (abhängig von Domäne)
% - Fazit: Vergleichbar zu anderen Techniken, nicht überlegen

% Subsubsection 3.2.3.3: "Nicht-Ersatz für gutes Lehren"
% \subsubsection{tDCS ersetzt nicht gutes Unterrichten}
\label{subsubsec:not_replacement}
% - tDCS ist ZUSATZ zu optimiertem Unterricht
% - Basissystem muss gut sein (sonst "garbage in, garbage out")
% - Kombination wichtig: tDCS + guter Trainer/Lehrer + strukturiertes Programm

\section{Praktikabilitäts-Fazit: Realistisches Szenario}
\label{sec:praktik_fazit}

% LITERATUR: Violante 2024 (kritikal), Cappon/Woodham für Machbarkeit

% \subsection{Machbarkeit nach Setting}
\label{subsec:machbarkeit_setting}

% Subsubsection 3.3.1.1: Élite-Training und spezialisierte Kontexte
% \subsubsection{Élite-Training (Sport, Musik, militärisch)}
\label{subsubsec:elite_training}
% - Hochmotivierte Populationen, kleinere Gruppen
% - Individuelle Betreuung möglich
% - Kosteneffizienz ggf. gegeben
% - Realistisch: 5-10 Jahre bis etablierte Praxis

% Subsubsection 3.3.1.2: Mainstream-Bildung (Schulen, Universitäten)
% \subsubsection{Mainstream-Bildung}
\label{subsubsec:mainstream_bildung}
% - Massive Skala (Millionen), Kostendruck
% - Heterogene Populationen (nicht alle profitieren)
% - Logistische Herausforderungen
% - Realistisch: 15-25 Jahre oder später für breite Integration

% Subsubsection 3.3.1.3: Home-Use für individuelle Lerner
% \subsubsection{Home-Use für Selbstoptimierung}
\label{subsubsec:homeuse_self}
% - Consumer-Geräte bereits verfügbar
% - Größtes Risiko: unsachgemäße Anwendung ohne Supervision
% - Dual-Use Problem: Medizinisch/Enhancement-Nutzung vermischt

%!TEX root = ../../main.tex
% ==============================================================================
% CHAPTER 4: ETHISCHE IMPLIKATIONEN
% ==============================================================================
\chapter{Ethische Implikationen und kritische Bewertung}
\label{chap:ethik}

% LITERATURVORSCHLÄGE FÜR KAPITEL:
% - Chatterjee (2013): "The ethics of neuroenhancement" [Standardwerk]
% - Brühl et al. (2019): "Neuroethical issues in cognitive enhancement" [37 Zit., balanced]
% - Dubljevic et al. (2016): "The is and ought of Ethics of Neuroenhancement" [Theory-Practice Gap]
% - Tubig & Racine (2024): "Cognitive Enhancement as Transformative Experience" [philosophisch, aktuell]
% - Cohen Kadosh (2015): Enhancement in Research/Rehabilitation/Skills Context
% - Violante et al. (2024): Kritische Bewertung Potential

\section{Zentrale ethische Dimensionen}
\label{sec:ethische_dimensionen}

% LITERATUR: Chatterjee 2013, Brühl 2019, Dubljevic 2016

% \subsection{Autonomie und informierte Einwilligung}
\label{subsec:autonomie_ic}

% LITERATUR: Chatterjee 2013, Tubig 2024, Cappon 2023

% Subsubsection 4.1.1.1: Freie Wahl in Machtstrukturen
% \subsubsection{Können Schüler/Trainees frei wählen?}
\label{subsubsec:free_choice}
% - Machtgefälle: Schüler vs. Lehrer, Athlet vs. Trainer, Mitarbeiter vs. Chef
% - Impliziter vs. expliziter Druck
% - Koerzion (Coercion): "Voluntary but nicht wirklich"?
% - Beispiel: "Nehmt tDCS an oder bekommt weniger Trainingsstunden"

% Subsubsection 4.1.1.2: Koerzion in kompetitiven Umgebungen
% \subsubsection{Koerzion in kompetitiven Umgebungen}
\label{subsubsec:koerzion_kompet}
% - Arms-Race-Dynamics: "Wenn andere tDCS nutzen, muss ich auch"
% - Pressure-Spirale in Leistungskulturen
% - "Collective Action Problem": Indiv. rational, system-irrational
% - Besonders problematisch: Junge Menschen (18-25)

% Subsubsection 4.1.1.3: Besonderheiten Minderjähriger
% \subsubsection{Spezielle Herausforderungen bei Kindern/Jugendlichen}
\label{subsubsec:minors_ic}
% - Hirnentwicklung bis ~25 Jahre
% - Nicht volle "Informed Consent" Capacity
% - Elterliche Einwilligung: kann diese Koerzion weiterleiten
% - International unterschiedlich reguliert

% Subsubsection 4.1.1.4: Anforderungen an echte Informed Consent
% \subsubsection{Standards für echte Informed Consent}
\label{subsubsec:ic_standards}
% - Transparente Information über Risiken/Nutzen/Unsicherheiten
% - Bewusste, freiwillige Entscheidung
% - Dokumentation und Opt-Out-Möglichkeit
% - Periodische Wiederbestätigung (nicht einmalig)

% \subsection{Gerechtigkeit und Zugang}
\label{subsec:gerechtigkeit}

% LITERATUR: Chatterjee 2013, Savulescu et al. 2011 (erwähnt in Cohen Kadosh 2015), Brühl 2019

% Subsubsection 4.1.2.1: "Enhancement-Divide" zwischen Populationen
% \subsubsection{Zugang und globale Gerechtigkeit}
\label{subsubsec:enhancement_divide}
% - Nur Vermögende haben Zugang (individuell oder institutionell)
% - Arme Länder/Communities ausgeschlossen
% - Verstärkt bestehende Ungleichheiten
% - Schulen mit Budget vs. ohne Budget

% Subsubsection 4.1.2.2: Wettbewerbsverzerrung
% \subsubsection{Wettbewerbsverzerrung und Fairness}
\label{subsubsec:unfairness}
% - Athleten/Musikschüler mit tDCS vs. ohne
% - Ist das ähnlich wie Doping?
% - Oder ähnlich wie teures Coaching/Equipment?
% - Unterschied zu genetischem Enhancement?

% Subsubsection 4.1.2.3: Gegenperspektive: Nicht zwangsläufig ungleichheitsverschärfend
% \subsubsection{Gegenperspektive: Potenzial für Egalitarismus}
\label{subsubsec:gegenperspektive_equal}
% - Cohen Kadosh (2015) argumentiert: Enhancement muss nicht Ungleichheit erhöhen
% - Wenn kostenlos/subventioniert verfügbar, könnte Chancen verbessern
% - Parallel-Investition in gute Schulen für alle notwendig

% Subsubsection 4.1.2.4: Regulatorische Lösungen
% \subsubsection{Mögliche regulatorische Lösungen}
\label{subsubsec:reg_solutions_access}
% - Öffentliche Finanzierung
% - Standardisierte Verfügbarkeit in staatlichen Systemen
% - Klare Kategorisierung: Therapie (finanziert) vs. Enhancement (User-Bezahlung)
% - Aber: Implementierungsrealität in verschiedenen Ländern unterschiedlich

% \subsection{Sicherheit und Langzeitfolgen}
\label{subsec:safety_longterm}

% LITERATUR: Brühl 2019, Cappon 2023, Woodham 2024 (nur 10 Wochen Daten)

% Subsubsection 4.1.3.1: Kurzfristige Sicherheit (gut dokumentiert)
% \subsubsection{Kurzfristige Nebenwirkungen (etabliert)}
\label{subsubsec:short_safety}
% - Mild, vorübergehend, reversibel
% - Klinische Praxis gut dokumentiert
% - Akzeptabler Risk-Benefit-Profile

% Subsubsection 4.1.3.2: Langfristige Sicherheit (unbekannt)
% \subsubsection{Langfristige Sicherheit (für Gesunde, unbekannt)}
\label{subsubsec:long_safety}
% - Brühl (2019): Forschungsbedarf für Lifetime Safety
% - Bisherige Studien: max. 5-10 Sessions
% - Woodham (2024): 10 Wochen Daten (längste bisher, aber immer noch kurz)
% - Fehlende: Prospektive Follow-Up nach 5-10 Jahren

% Subsubsection 4.1.3.3: Hirnentwicklung und Neuroplastizität
% \subsubsection{Besondere Bedenken bei Hirnentwicklung (<25 Jahre)}
\label{subsubsec:brain_development}
% - Präfrontalkortex entwickelt sich bis ~25 Jahre
% - Adoleszente Neuroplastizität anders als Erwachsene
% - Könnte wiederholte Stimulation Entwicklung negativ beeinflussen? (Spekulativ)
% - Forschungsbedarf: Spezifische Sicherheitsstudien in Adoleszenten

% Subsubsection 4.1.3.4: Mögliche unerwartete Langzeitkonsequenzen
% \subsubsection{Spekulation: Mögliche unerwartete Konsequenzen}
\label{subsubsec:unexpected_effects}
% - Neuroplastisches "Over-Training"? (Ähnlich wie muskuläre Überanstrengung)
% - Abhängigkeitsentwicklung? (psychologisch, nicht physiologisch)
% - Shift zu "Technik-abhängigen" Lernmechanismen?
% - ABER: Spekulativ, keine Evidenz - aber erfordert Forschung

% \subsection{Authentizität, Leistung und menschliche Exzellenz}
\label{subsec:authenticity}

% LITERATUR: Chatterjee 2013, Cohen Kadosh 2015, Tubig 2024

% Subsubsection 4.1.4.1: Die philosophische Frage
% \subsubsection{Wer hat geleistet - ich oder das Gerät?}
\label{subsubsec:agency_question}
% - Zentrale Sorge: Ist tDCS-enhancte Leistung "wirklich" meine Leistung?
% - Ähnlich wie: Ist mit Taschenrechner gelöste Matheaufgabe meine Leistung?
% - Oder mit Spellchecker geschriebener Essay?
% - Philosophisch umstritten

% Subsubsection 4.1.4.2: Argumente PRO Enhancement (auch mit tDCS)
% \subsubsection{Argumente FÜR Enhancement als authentisch}
\label{subsubsec:pro_enhancement}
% - Cohen Kadosh (2015): Enhancement ähnlich wie Schlafhygiene, richtige Ernährung, Coaching
% - Person trainiert trotzdem selbst, tDCS ist nur Unterstützung
% - Nicht "medikalisiert" wie Modafinil während Aufgabe
% - Vereinbar mit Authentizität (wenn Information transparent)

% Subsubsection 4.1.4.3: Argumente GEGEN Enhancement
% \subsubsection{Argumente GEGEN Enhancement (Sorgen)}
\label{subsubsec:contra_enhancement}
% - Grundlegende Verschiebung: von "Mensch leistet" zu "Mensch + Technik leistet"
% - Medikalisierung von Leistung (medicalization of performance)
% - Fokusverschiebung: von intrinsischen Fähigkeiten zu technologischer Assistenz
% - "Race to the Bottom": Wettbewerb wird Geräte-abhängig statt Mensch-abhängig

% Subsubsection 4.1.4.4: Nuancierte Position
% \subsubsection{Nuancierte Position}
\label{subsubsec:nuanced_position_auth}
% - Unterscheiden: Type von Enhancement (während Trainig vs. während Test)
% - Unterscheiden: Kontext (Elite-Sport hat andere Standards als Schule)
% - Transparenz kritisch: Müssen andere wissen, dass tDCS verwendet wurde?
% - Balance: Unterstützung ermöglichen, aber nicht als "Shortcut"

% \subsection{Charakter, Anstrengung und menschliche Tugenden}
\label{subsec:character}

% LITERATUR: Chatterjee 2013, Brühl 2019, philosophische Editorials

% Subsubsection 4.1.5.1: Traditionelle Ansicht
% \subsubsection{Traditionelle Perspektive: Anstrengung bildet Charakter}
\label{subsubsec:trad_character}
% - Klassisches Verständnis: Schwierigkeiten überwinden = Charakter-Bildung
% - Resilienz, Durchhaltevermögen, Grit entwickeln sich durch Kampf
% - Enhancement könnte diese Entwicklung unterminieren
% - Risiko: "Zu bequeme" Leistung

% Subsubsection 4.1.5.2: Pragmatische Perspektive
% \subsubsection{Pragmatische Perspektive: Effiziente Lernmethoden erwünscht}
\label{subsubsec:pragmatic_efficiency}
% - Wenn tDCS Lernen beschleunigt: Mehr Zeit für andere Fähigkeiten
% - Effizienzgewinn gesellschaftlich wertvoll
% - Nicht alle Anstrengung notwendig (distingt von sinnvoller Herausforderung)
% - Balance: Minimum-Anstrengung für Resilienz, dann tDCS zur Optimierung

% Subsubsection 4.1.5.3: Balanceakt
% \subsubsection{Balanceakt: Unterstützung ohne Atrophie}
\label{subsubsec:balance_act}
% - Notwendig: Kontinuierliche Herausforderung
% - tDCS als "Beschleuniger" nicht als "Ersatz für Anstrengung"
% - Gefahr: Komplett auf tDCS verlassen = Skill-Atrophie
% - Gesellschaftlich: Wenn alle tDCS nutzen, "resets" der relative Vorteil

\section{Ethische Regulierung und Governance}
\label{sec:governance}

% LITERATURVORSCHLÄGE: Dubljevic 2016, Cappon 2023 (praktische Implementation), neurocare academy

% \subsection{Status quo: Regulatorische Grauzonen}
\label{subsec:status_quo}

% LITERATUR: FDA Guidelines (implicit), IFCN Guidelines, Commercialization risks

% Subsubsection 4.2.1.1: Unterschied Therapie vs. Enhancement
% \subsubsection{Regulatorische Unterscheidung: Therapie vs. Enhancement}
\label{subsubsec:therapy_vs_enhancement}
% - Therapie (FDA-approved): Depression (rTMS), Migräne, etc.
% - Enhancement (Gray Zone): Kognitives Lernen, Skill Acquisition
% - Problem: Gleiche Technologie, sehr unterschiedliche Regulierung
% - Inkonsistenz: Motorisches Lernen "Enhancement", aber Schlaganfall-Rehab "Therapie" (neuro-Restauration)

% Subsubsection 4.2.1.2: Consumer-Geräte und DIY-Kultur
% \subsubsection{Consumer-Geräte und unkontrollierte Anwendung}
\label{subsubsec:consumer_diys}
% - Amazon/eBay: tDCS-Kits kaufbar (~$300)
% - Internet: DIY-Konstruktions-Tutorials
% - Maximales Risiko: Unsachgemäße Elektrodenplatzierung, falsche Parameter
% - Minimales Überwachung: Keine medizinische Aufsicht

% Subsubsection 4.2.1.3: Internationale Unterschiede
% \subsubsection{Internationale Unterschiede in Regulierung}
\label{subsubsec:intl_differences}
% - USA: FDA-Zulassung für spezifische therapeutische Indikationen
% - Europa: Variabel (Deutschland, Schweiz, UK unterschiedlich)
% - Andere Länder: Oft wenig Regulation
% - Problem: "Regulatory Arbitrage" (Leute gehen zur lockereren Regulierung)

% \subsection{Empfohlene Regulierungsrahmen}
\label{subsec:empf_reg_frameworks}

% LITERATUR: Dubljevic 2016 (ethisch), IFCN Guidelines (praktisch), Cappon 2023 (Implementation)

% Subsubsection 4.2.2.1: Prä-Marktzulassung Standards
% \subsubsection{Evidenzstandards vor Markteinführung}
\label{subsubsec:premarket_standards}
% - Sicherheit (Kurzzeit + längerfristig): Minimal 12 Wochen Studien
% - Effektivität: Randomized Controlled Trials
% - Risiko-Benefit-Analyse
% - Unterschied: Therapie (höherer Standard) vs. Enhancement (möglicherweise niedrigerer)

% Subsubsection 4.2.2.2: Informed Consent Standards
% \subsubsection{Strikte Informed Consent Anforderungen}
\label{subsubsec:ic_standards_reg}
% - Schriftlich, nicht mündlich
% - Laigerechte Sprache (nicht nur medizinische Fachbegriffe)
% - Transparente Angabe von Unsicherheiten ("nicht vollständig untersucht")
% - Spezielle Standards für Minderjährige
% - Opt-Out jederzeit möglich

% Subsubsection 4.2.2.3: Fairness und Zugangs-Guidelines
% \subsubsection{Fairness-Guidelines und Zugang}
\label{subsubsec:fairness_guidelines}
% - Definition: Wann ist Enhancement "fair" vs. "unfair"?
% - Sport: Ähnliche Regelwerk wie andere Performance-Enhancers?
% - Bildung: Verboten in Prüfungen? Oder in Lernphase erlaubt?
% - Gleichzeitiger Zugang: Alle sollten Möglichkeit haben (falls legal)

% Subsubsection 4.2.2.4: Transparenz und Dokumentation
% \subsubsection{Transparenz und Dokumentation}
\label{subsubsec:transparency}
% - Wenn tDCS in Setting (Schule, Club) verwendet: Dokumentation nötig
% - Wer weiß davon? (Athletes, anderen Schülern, Regulierungsbehörden)
% - Verpflichtete Berichterstellung von Nebenwirkungen
% - Public Databases für Enhancement-Use (ähnlich wie Pharmakovigilanz)

% Subsubsection 4.2.2.5: Langzeitüberwachung
% \subsubsection{Registries und Langzeitüberwachung}
\label{subsubsec:registries}
% - Enhancement-Registries: Wer verwendet, wie oft, mit welchen Effekten
% - Cohort-Studies: Follow-Up über Jahre
% - Safety Monitoring: kontinuierliche Überwachung
% - Datenschutz: Balance zwischen Überwachung und Privacy

% \subsection{Ethische Governance-Strukturen}
\label{subsec:governance_structures}

% LITERATUR: Dubljevic 2016 (Stakeholder-Integration), Ethics Committee Guidelines

% Subsubsection 4.2.3.1: Interdisziplinäre Kommissionen
% \subsubsection{Interdisziplinäre Ethik-Kommissionen}
\label{subsubsec:interdisc_commissions}
% - Zusammensetzung: Neurowissenschaftler, Ethiker, Pädagogen, Juristen, Patienten-Vertreter
% - Aufgabe: Policy-Entwicklung, Fall-by-Case Review, Konfliktresolution
% - Transparent: Öffentliche Beratungen

% Subsubsection 4.2.3.2: Stakeholder-Partizipation
% \subsubsection{Integration von Stakeholder-Perspektiven}
\label{subsubsec:stakeholder_participation}
% - Dubljevic (2016): Derzeitige Debatte vernachlässigt Stakeholder
% - Schüler, Eltern, Lehrer, Athleten: Ihr Input kritisch
% - Nicht-Experten-Perspektiven wertvoll
% - Partizipative Entscheidungsfindung (nicht Top-Down)

% Subsubsection 4.2.3.3: "Precautionary Principle" vs. Innovation
% \subsubsection{"Precautionary Principle" balanciert mit Innovation}
\label{subsubsec:precautionary}
% - Zu konservativ: Blockiert nützliche Innovationen
% - Zu liberal: Akzeptiert ungeprüfte Techniken
% - Balance: Moderat-Precautionary Approach
%   - Zulassen in Forschung (mit Oversight)
%   - Vorsicht bei breiter Publikums-Anwendung
%   - Längerfristig: Flexibilität mit neuen Daten

\section{Normative Position und ethische Schlussfolgerungen}
\label{sec:normative}

% LITERATUR: Chatterjee 2013, Brühl 2019, Cohen Kadosh 2015, Tubig 2024

% \subsection{Unter welchen Bedingungen ethisch vertretbar?}
\label{subsec:when_justified}

% Subsubsection 4.3.1.1: Notwendige Bedingungen
% \subsubsection{Notwendige und hinreichende Bedingungen für Ethische Vertretbarkeit}
\label{subsubsec:necessary_conditions}
% 1. Freiwilligkeit (keine Koerzion)
% 2. Informierte Zustimmung (transparent, verständlich)
% 3. Zumutbares Sicherheitsprofil (Risik-Benefit positiv)
% 4. Gerechtigkeit (äquitable Zugang oder Fairness-Kriterium)
% 5. Transparenz (anderen bekannt, wenn relevant)
% 6. Professionelle Aufsicht (nicht DIY ohne Kontrolle)
% 7. Reversibilität (kann man stoppen ohne Schaden?)

% Subsubsection 4.3.1.2: Zusätzliche Bedingungen für Kinder/Jugendliche
% \subsubsection{Zusätzliche Standards für Minderjährige}
\label{subsubsec:minor_conditions}
% - Elterliche Zustimmung + unabhängige ethische Prüfung
% - Besondere Sicherheitsforschung für Adoleszenten
% - Eindeutige medizinische Indikation (nicht bloß Enhancement)
% - Regelmäßige neuro-psychologische Monitoring

% \subsection{Kritische Kontexte und deren Bewertung}
\label{subsec:critical_contexts}

% Subsubsection 4.3.2.1: Eliteschule / Hochleistungssport
% \subsubsection{Eliteschule und Hochleistungssport}
\label{subsubsec:elite_sport}
% - Hohes Koerzions-Risiko ("Druck, um konkurrenzfähig zu bleiben")
% - Junge Menschen (17-25) besonders vulnerabel
% - Urteil: ETHISCH FRAGLICH, nicht empfohlen ohne strikte Safeguards
% - Falls erlaubt: Nur mit unabhängigem IC, Anti-Koerzions-Monitoring

% Subsubsection 4.3.2.2: Hochschule unter Exam-Stress
% \subsubsection{Hochschule unter Prüfungs-Stress}
\label{subsubsec:university_stress}
% - Mittleres Koerzions-Risiko
% - Autonomere Individuen (20-30 Jahre)
% - Urteil: AMBIVALENT
%   - IF volle IC + keine institution. Unterstützung = eher akzeptabel
%   - IF Norm wird "alle nutzen tDCS" = problematisch
% - Empfehlung: Parallel-Ansätze zur Stressreduktion (bessere Prüfungsformate, mental health support)

% Subsubsection 4.3.2.3: Rehabilitations-Setting (post-Schlaganfall)
% \subsubsection{Rehabilitation nach Schlaganfall}
\label{subsubsec:rehab_stroke}
% - THERAPEUTISCH, nicht Enhancement
% - Andere ethische Kategorie (Wiederherstellung von Funktion)
% - Urteil: ETHISCH GERECHTFERTIGT (wenn sicher und effektiv)
% - Weniger bedenken als Enhancement

% Subsubsection 4.3.2.4: Kinder/Jugendliche in allgemeiner Schule
% \subsubsection{Kinder und Jugendliche in allgemeiner Schule}
\label{subsubsec:children_school}
% - Höchstes Koerzions-Risiko + Entwicklungs-Bedenken
% - Urteil: NICHT EMPFOHLEN für Enhancement
% - Exception: Therapeutische Indikation (z.B. ADHS, Dyslexia)
% - Grund: Zu viele Unsicherheiten, zu vulnerable Population

% \subsection{Offene ethische Fragen}
\label{subsec:open_ethical_questions}

% Subsubsection 4.3.3.1: Philosophische Unterscheidung zu anderen Enhancements
% \subsubsection{Philosophische Unterscheidungen}
\label{subsubsec:phil_distinctions}
% - Warum ist tDCS ethisch anders als Modafinil?
% - Warum anders als mentales Training, Coaching, Mentaltraining?
% - Warum anders als teures Equipment im Sport?
% - Oder ist es gar NICHT so unterschiedlich?

% Subsubsection 4.3.3.2: Gesellschaftliche Auswirkungen bei breiter Verbreitung
% \subsubsection{Gesellschaftliche Szenarien}
\label{subsubsec:societal_scenarios}
% - Szenario A: "Spiral to the Bottom" (alle müssen tDCS nutzen zum Mithalten)
% - Szenario B: "Enhancement as Luxury" (nur für Wohlhabende, verstärkt Ungleichheit)
% - Szenario C: "Public Good" (subventioniert, alle haben Zugang, neue Normal)
% - Unklar, welches Szenario realistisch ist

% Subsubsection 4.3.3.3: Internationale und kulturelle Unterschiede
% \subsubsection{Kulturelle Unterschiede}
\label{subsubsec:cultural_diffs}
% - Unterschiedliche Ethik-Standards weltweit
% - "Enhancement" ist kulturell-konstruiert (nicht universal)
% - Problem: Globale Technologie, lokal unterschiedliche Ethik
% - Notwendigkeit: Interkultureller Dialog, aber Respekt für Variabilität

%!TEX root = ../../main.tex
% ==============================================================================
% CHAPTER 5: ZUSAMMENFASSUNG UND FAZIT
% ==============================================================================
\chapter{Zusammenfassung und Fazit}
\label{chap:fazit}

% LITERATURVORSCHLÄGE FÜR KAPITEL:
% - Violante et al. (2024): "Can neurotechnology revolutionize cognitive enhancement?" [überblick, kritisch]
% - Brühl et al. (2019): Balanced Assessment [Sicherheit, Ethik]
% - Cappon/Woodham: praktische Perspektive

\section{Beantwortung der Forschungsfrage}
\label{sec:beantwortung}

% LITERATUR: Zu all sections passend

% \subsection{These 1: Praktikabilität - JA, MIT EINSCHRÄNKUNGEN}
\label{subsec:thesis_praktikabilität}

% Subsubsection 5.1.1.1: Technische Machbarkeit
% \subsubsection{Technische Machbarkeit}
\label{subsubsec:tech_feasible}
% - Evidenz: Cappon 2023, Woodham 2024
% - Home-based, scaleable, niedrige Nebenwirkungen
% - Laien-Administration möglich mit Training

% Subsubsection 5.1.1.2: Wirksame Fertigkeits-Beschleunigung in bestimmten Domänen
% \subsubsection{Effektivität}
\label{subsubsec:effectiveness_summary}
% - Motorisches Lernen: ~15-30% Beschleunigung [Reis, Akkad]
% - Akademisch: SMD ~0.3-0.7 [Simonsmeier]
% - Arbeitsgedächtnis: SMD ~0.3-0.5 [Senkowski]
% - ABER: 60% Non-Responder, Variabilität hoch

% Subsubsection 5.1.1.3: Haupt-Limitationen
% \subsubsection{Praktische Limitationen}
\label{subsubsec:limitations_summary}
% - Inter-individuelle Variabilität: 60% Non-Responder
% - Kosten für Skalierung
% - Personalisierungs-Bedarf
% - Nicht "Wundermittel" - ersetzt nicht gutes Lehren

% Subsubsection 5.1.1.4: Realistisches Szenario
% \subsubsection{Realistischer Zeitrahmen und Szenarios}
\label{subsubsec:timeframe}
% - Élite-Kontexte (Sport, Musik, Militär): 5-10 Jahre bis etabliert
% - Mainstream-Bildung: 15-25+ Jahre (wenn je)
% - Continuous home-use: Schon möglich, aber minimal überwacht

% \subsection{These 2: Ethische Implikationen - KOMPLEX, ABER VERWALTBAR}
\label{subsec:thesis_ethik}

% Subsubsection 5.1.2.1: Nicht fundamental unethisch
% \subsubsection{Nicht grundsätzlich unethisch}
\label{subsubsec:not_fundamentally_unethical}
% - Prinzipiell vertretbar unter bestimmten Bedingungen
% - Ähnlich zu anderen akzeptierten Enhancements
% - Aber: Stark kontext-abhängig

% Subsubsection 5.1.2.2: Substanzielle Bedenken EXISTIEREN
% \subsubsection{Substantielle ethische Herausforderungen}
\label{subsubsec:substantive_concerns}
% - Koerzion in kompetitiven Settings
% - Gerechtigkeit und Zugang
% - Langzeit-Sicherheit unbekannt
% - Authentizität-Fragen

% Subsubsection 5.1.2.3: Zentrale Anforderungen
% \subsubsection{Zentrale Anforderungen für Ethische Vertretbarkeit}
\label{subsubsec:central_requirements}
% 1. Freiwilligkeit (Anti-Koerzion)
% 2. Informierte Zustimmung (transparent, verstanden)
% 3. Sicherheit (akzeptables Risik-Benefit)
% 4. Gerechtigkeit (äquitable Zugang oder klare Fairness-Regeln)
% 5. Transparenz (Offenlegung bei relevant)

% Subsubsection 5.1.2.4: Hohe Risiken in spezifischen Kontexten
% \subsubsection{Kontexte mit hohem Risiko}
\label{subsubsec:high_risk_contexts}
% - Eliteschule + junge Menschen = NICHT EMPFOHLEN
% - Mainstream-Schule + Minderjährige = NICHT EMPFOHLEN
% - Rehabilitation = ANDERS (therapeutisch, eher gerechtfertigt)
% - Informed Erwachsene im Eigeninteresse = ambivalent, aber ggf. akzeptabel

% Subsubsection 5.1.2.5: Regulierung essentiell
% \subsubsection{Regulierung ist essentiell}
\label{subsubsec:regulation_essential}
% - Nicht nur Marktlösung
% - Nicht nur Selbstregulation
% - Governance-Strukturen notwendig

\section{Synthese und Übergeordnete Erkenntnisse}
\label{sec:synthesis}

% LITERATUR: Violante 2024, Brühl 2019

% \subsection{Realistische Einschätzung der Technologie}
\label{subsec:realistic_assessment}

% Subsubsection 5.2.1.1: Weder "Wundermittel" noch "Insignifikant"
% \subsubsection{Mittlere Position zwischen Hype und Skeptizismus}
\label{subsubsec:middle_position}
% - NICHT: "Hirnstimulation wird Bildung/Training revolutionieren"
% - NICHT: "Hirnstimulation hat keine praktischen Anwendungen"
% - EHER: "Nützliches Zusatz-Werkzeug für spezifische Kontexte"
% - Violante (2024): "modest evidence" ist fairere Beschreibung

% Subsubsection 5.2.1.2: Größeres Problem ist System, nicht Technologie
% \subsubsection{Systemische Fragen wichtiger als Technologie-Fragen}
\label{subsubsec:system_questions}
% - Warum Leistungsdruck so hoch?
% - Warum so viele Schüler "brauchen" Enhancement?
% - Sind schlecht designte Schulen/Trainings das Kernproblem?
% - Parallel: Systeme verbessern + Technologie verwenden (nicht nur letzteres)

% \subsection{Unterschied zwischen Enhancement-Anwendung und Enhancement-Notwendigkeit}
\label{subsec:use_vs_need}

% Subsubsection 5.2.2.1: Technologie ermöglicht, aber macht nicht notwendig
% \subsubsection{Technology Push vs. Demand Pull}
\label{subsubsec:tech_push_pull}
% - Technologie existiert (Tech-Push)
% - Gesellschaftlicher Druck nach Enhancement (Demand-Pull)
% - Beides zusammen = rapid adoption risk
% - Lösungsansatz: Demand reduzieren (bessere Systeme), nicht nur Supply regulieren

\section{Empfehlungen für die Zukunft}
\label{sec:empfehlungen}

% LITERATUR: Cappon 2023, Woodham 2024 (praktisch), Brühl 2019, Dubljevic 2016 (ethisch)

% \subsection{Wissenschaftliche Empfehlungen}
\label{subsec:wiss_empf}

% Subsubsection 5.3.1.1: Langzeitstudien
% \subsubsection{Prospektive Langzeitstudien zur Sicherheit}
\label{subsubsec:long_term_studies}
% - Mindestens 2-5 Jahre Follow-Up
% - Fokus: Gesunde Population, repeated use
% - Spezial-Kohorten: Adoleszente (<25 Jahre), ältere Erwachsene
% - Neuro-psychologische Assessments, neuroimaging

% Subsubsection 5.3.1.2: Biomarker-Forschung
% \subsubsection{Responder-Biomarker identifizieren}
\label{subsubsec:biomarker_research}
% - Wer profitiert? Genetik, fMRI, EEG Baseline?
% - Prädiktive Modelle entwickeln
% - Validierung in unabhängigen Kohorten
% - Ziel: Personalisierte Protokolle

% Subsubsection 5.3.1.3: Transfer-Effekte
% \subsubsection{Transfer- und Generalisierungs-Forschung}
\label{subsubsec:transfer_research}
% - Wechselt Lernen zu neuem Kontext?
% - Vergleich: Lab vs. Realwelt
% - Mechanismen verstehen

% Subsubsection 5.3.1.4: Vergleichsstudien
% \subsubsection{Benchmarking gegen etablierte Methoden}
\label{subsubsec:benchmarking_research}
% - tDCS vs. Spaced Repetition
% - tDCS vs. Elaboration
% - tDCS + Methode vs. Methode allein?

% \subsection{Praktische Empfehlungen}
\label{subsec:prakt_empf_final}

% Subsubsection 5.3.2.1: Pilot-Programme in kontrollierten Settings
% \subsubsection{Strukturierte Pilotierungen}
\label{subsubsec:structured_pilots}
% - Elite-Sport Teams (mit IC, Monitoring)
% - Spezialmusikschulen
% - Militär-Trainingsprogramme
% - Mit ethischer Überwachung
% - Published Results (nicht hidden failures)

% Subsubsection 5.3.2.2: Home-Based mit Professionellem Support
% \subsubsection{Home-Based Modelle mit Fernüberwachung}
\label{subsubsec:homebased_models}
% - Cappon/Woodham Modelle als Template
% - Trainingsprogramme für Anwender
% - Remote-Ärzte/Fachpersonen für Support
% - Clear Safety Protocols

% Subsubsection 5.3.2.3: Integration mit optimierten Lehrmethoden
% \subsubsection{Nicht Technologie-Zentriert}
\label{subsubsec:not_techno_centric}
% - tDCS ist ADD-ON, nicht Fokus
% - Parallel: Pädagogische Optimierung
%   - Besserer Unterricht, bessere Trainer
%   - Optimale Lernumgebung
%   - Mentale Health Support
% - tDCS nur wenn fundamentale Strukturen gut sind

% \subsection{Regulatorische Empfehlungen}
\label{subsec:reg_empf_final}

% Subsubsection 5.3.3.1: Nationale Guidelines entwickeln
% \subsubsection{Evidenzbasierte nationale Guidelines}
\label{subsubsec:national_guidelines}
% - Pro Land: Nationale Policy
% - Basierend auf lokal akzeptierten ethischen Standards
% - ABER: Harmonisierung wo möglich (z.B. EU)

% Subsubsection 5.3.3.2: Klare Therapie/Enhancement Abgrenzung
% \subsubsection{Regulatorische Unterscheidung verfeinern}
\label{subsubsec:therapy_enhancement_distinction}
% - Therapie: Wiederherstellung von Funktion (höherer Regulierungs-Standard)
% - Enhancement: Verbesserung von normal-range Funktionen (niedrigerer Standard)
% - ABER: Grauzone groß (z.B. motorische Rehabilitation vs. Athletisches Training)
% - Einzelfall-Betrachtung notwendig

% Subsubsection 5.3.3.3: Strikte Informed Consent in Enhancement-Anwendungen
% \subsubsection{IC-Standards für Enhancement}
\label{subsubsec:ic_enhancement_standards}
% - Schriftlich, in Laiensprache
% - Unsicherheiten transparent
% - Speziell für Minderjährige: unabhängige ethische Prüfung
% - Dokumentation und regelmäßige Bestätigung

% Subsubsection 5.3.3.4: Monitoring und Registries
% \subsubsection{Tracking und Langzeit-Überwachung}
\label{subsubsec:monitoring_final}
% - Registry für Enhancement-User (wer, wie oft, welche Effekte)
% - Safety-Reporting (adverse events)
% - Datenbank für Forschung + Regulierung
% - Datenschutz: Privacy gewährleisten

% \subsection{Ethische Empfehlungen}
\label{subsec:eth_empf_final}

% Subsubsection 5.3.4.1: Gesellschaftliche Debatte führen
% \subsubsection{Offene gesellschaftliche Diskussion}
\label{subsubsec:public_debate}
% - Nicht nur Expert*innen, auch Laien einbeziehen
% - Media literacy: Hype vs. Realität
% - Stakeholder-Partizipation (Schüler, Eltern, Trainer)
% - Interdisziplinäre Dialoge

% Subsubsection 5.3.4.2: Langzeit-Folgen-Monitoring bei Early Adopters
% \subsubsection{Cohort-Tracking von Early Adopters}
\label{subsubsec:cohort_tracking}
% - Erste Menschen, die regelmäßig tDCS nutzen: folgen
% - 5-10 Jahre Neuro-psychologische Assessments
% - Identifikation unerwarteter Langzeiteffekte

% Subsubsection 5.3.4.3: "Precautionary Principle" als Leitlinie
% \subsubsection{Moderates Precautionary Approach}
\label{subsubsec:precaut_leitlinie}
% - Nicht blockierend, aber vorsichtig
% - Forschung erlauben (mit Oversight)
% - Breite Publikums-Nutzung: Vorsicht bis mehr Daten
% - Flexibilität mit neuen Evidenzen

% Subsubsection 5.3.4.4: Parallel: Systemischen Druck reduzieren
% \subsubsection{Systemic Change parallel zu Technologie}
\label{subsubsec:systemic_change}
% - Nicht nur: "Wie use tDCS responsibel?"
% - Auch: "Warum ist Enhancement nötig?"
% - Investieren in:
%   - Bessere Schulen, bessere Lehrer
%   - Weniger Prüfungsdruck
%   - Mental Health Support
%   - Bessere Work-Life-Balance in Elite-Training
% - Technologie + Systemische Reformen zusammen

\section{Abschließende Bewertung und Forschungsfrage-Antwort}
\label{sec:abschließend}

% LITERATUR: Zusammenfassung aus allen Kapiteln

% \subsection{Zentrale Erkenntnis}
\label{subsec:zentraleerkn}

% Text placeholder für Zusammenfassung und abschließende These

% \subsection{Zukunftsausblick}
\label{subsec:ausblick}

% Subsubsection 5.4.2.1: Wahrscheinliche Entwicklung
% \subsubsection{Realistische Zukunfts-Szenarien (nächste 10-20 Jahre)}
\label{subsubsec:future_scenarios}
% - Szenario 1: "Nischenproduzent" (nur Elite, Therapie, Forschung)
% - Szenario 2: "Unter-regulated Mass Market" (DIY, Kommerzialisierung, Risiken)
% - Szenario 3: "Regulated Enhancement Tool" (Guidelines, Governance, limitierte Nutzung)
% - Welcher Szenario am wahrscheinlichsten?

% Subsubsection 5.4.2.2: Kritische Faktoren
% \subsubsection{Kritische Faktoren für Outcome}
\label{subsubsec:critical_factors}
% - Forschungs-Fortschritt (z.B. bessere Biomarker -> Personalisierung)
% - Regulatorische Entscheidungen (z.B. FDA policy)
% - Gesellschaftliche Werte (z.B. Akzeptanz von Enhancement)
% - Ökonomische Faktoren (Kosten-Reduktion oder nicht?)

% Subsubsection 5.4.2.3: Finale Empfehlung
% \subsubsection{Finale Position}
\label{subsubsec:final_position}
% - Dafür: Responsibel nutzen wo Evidenz existiert
% - Dagegen: Massiver Hype, unkontrollierte Nutzung, vulnerable Populationen
% - Neurale-Balance: Innovation ermöglichen, aber mit Guardrails


%TC:ignore
\iffalse
	\chapter{First Chapter}

	Das \gls{abstract} beschreibt in wenigen Sätzen die Zielsetzung und das Ergebnis der Ausarbeitung. Das Abstract muss sich vollständig auf der Titelseite befinden. Die Zeichensatzformatierung wird in einem eigenen Absatz beschrieben  Das Abstract soll es den Lesern:innen ermöglichen, innerhalb von wenigen Augenblicken zu erfassen, welcher Inhalt hinter der Überschrift steckt und ob das Thema, aus Sicht der Leser:innen, zur weiteren Bearbeitung lohnt. Das Abstract ist keine verbale Beschreibung des Inhaltsverzeichnisses, sondern gibt kurz und knapp z.B. die Zielsetzung (z.B. Hypothese), die eingesetzten Methoden und die erzielten Ergebnisse / Erkenntnisse bekannt. Weitere Hinweise finden Sie außerdem im Vorlesungsskript.

	Beispielverweis auf Quelle \cite{ahrensAbschlussarbeiten}.

	Test-Acronym: \acrshort{gcd}.

	\section{First Section}

	% Beispieltabelle
	\begin{table}[h]
		\centering
		\begin{tabular}{|c|c|}
			\hline
			Spalte 1 & Spalte 2 \\
			\hline
			Inhalt 1 & Inhalt 2 \\
			\hline
		\end{tabular}
		\caption{Beispiel-Tabelle}
		\label{tab:example}
	\end{table}

	\chapter{Second Chapter}

	\section{Another Section}

	% Beispielabbildung
	\begin{figure}[h]
		\centering
		\includegraphics[width=0.5\textwidth]{HSRTReport/Assets/Images/METI.png}
		\caption{MeTI-Logo}
		\label{fig:meti}
	\end{figure}

	% Alles hier wird von TexCount ignoriert.
	%TC:endignore
\fi
